\documentclass[11pt]{article}
\usepackage[T1]{fontenc}
\usepackage[margin=1.25in]{geometry}
\usepackage[expansion=false]{microtype}
\usepackage{amsmath,amssymb}
\usepackage{booktabs,tabularx,array}
\usepackage{enumitem}
\usepackage{hyperref}
\emergencystretch=3em
\setlength{\parskip}{0.55em}
\setlength{\parindent}{0pt}

\title{The Consensus Engine}
\author{Flyxion\\\small Independent Researcher}
\date{30 September 2026}

\begin{document}
\maketitle

\begin{abstract}
\noindent
The Consensus Engine stores a sentence if and only if more than half of the respondents who answered say that it is true. It stores nothing else, which its Office of Settled Questions regards as its chief economy. This report follows the Engine through its first fourteen years and its ten standing committees, and records what each committee decided and what the decisions left behind. The Engine, the Office, the committees, the sentences and every number other than plain arithmetic are invented, and nothing here describes any real system or person. The failures that accumulate are nonetheless of four familiar kinds, concerning units, scope, negation and scalar status, and the last section restates them in the terms of the monograph \emph{Adversaria: A Specification for a Public Claim Graph}. The thesis, which the Engine confirmed at its own expense, is that a system which stores only what a majority says is true will eventually store nothing it can explain.
\end{abstract}

\section{The principle of storage}

The founding document of the Engine is one paragraph long and is reproduced here in full. ``A sentence is entered in the Engine if more than half of those who answer agree that it is true. It is then Settled. A sentence that fails this test is not entered. Nothing else is entered.''

The merits of the arrangement were stated at the opening ceremony. The record of each sentence is a string and a tally, two fields, so that the Engine has no schema to maintain. There is no need to store reasons, since a sentence that most people believe has, by that fact, reasons enough. There is no need to store the limits of a sentence, since a sentence Settled is Settled. And there is no need to store falsehoods at all, since they are precisely what is not Settled. Visitors to the Office were told that the Engine had the smallest footprint of any knowledge system in existence, which was true, and that it never contradicted itself, which was not checked.

The Office's glossary defines ``Settled'' as follows: that someone has accepted the sentence; that the sentence follows from what is already known; and that most respondents agree. A poll on which of the three senses was intended closed at 34, 33 and 33, and the glossary entry was amended to read ``Settled: see above.''

The work of the Engine is governed by the committees in the following table. Each committee was constituted by a vote, as were the rules for constituting them.

\begin{center}
\small
\begin{tabularx}{\textwidth}{@{}l X X@{}}
\toprule
committee & mandate & later history\\
\midrule
Committee on Agreement & decides who counts as a respondent & settled that a respondent is whatever answers\\
Subcommittee on Quorum & decides how many must be present & quorate when a majority of those present attend\\
Committee on Units & decides what a sentence is & settled that a sentence is whatever fits the field\\
Committee on Scope & proposed that sentences say where they hold & abolished in Year 2 (Section~\ref{sec:scope})\\
Joint Committee on Negation & decides what is done with sentences that are not liked & issued one ruling (Section~\ref{sec:negation})\\
Committee on Contradiction & decides what is done with pairs of Settled sentences that disagree & its existence was Settled at 52 percent and its non-existence at 54\\
Panel on Weights & decides how Settled sentences are ranked & see Section~\ref{sec:scalar}\\
Committee on Reasons & records why sentences are Settled & never quorate (Section~\ref{sec:reasons})\\
Office of Appeals & hears appeals & hears only appeals that a majority agrees are appealable\\
Committee for the Review of Committees & reviews the committees & reviewed itself; the review is Settled\\
\bottomrule
\end{tabularx}
\end{center}

\section{The unit}\label{sec:units}

The Committee on Units met first and met briefly. Its report is three lines long: a sentence is a string; a string is whatever can be typed into the field; two strings are two sentences. The Committee observed that it had thereby settled the unit of knowledge without consulting anyone who knew anything, which it regarded as a saving.

The consequences appeared within the first year. The sentence ``The Tuesday market is open'' was entered with a tally of 68 of 100. The sentence ``On Tuesdays the market is open'' was entered separately, with 61. Neither entry knew of the other. A paraphrase of a Settled sentence, it emerged, was Settled about 40 percent of the time, and the Office attributed this to carelessness in paraphrasing and not to the design.

A second consequence concerned aggregates. A respondent submitted the sentence ``Soup is, on the whole, liked,'' together with the observation, which the Engine did not store, that nine of ten households had rated a certain soup $+1$ and the tenth had rated it $-5$. The total was $9-5=4$, which is positive, and the sentence was Settled at 71 percent. A later visitor asked whether the tenth household liked the soup. The Engine replied ``Settled: yes.'' The Committee on Units was asked to comment and pointed out that the question concerned a different sentence. The visitor pointed out that the Engine had no means of knowing that.

\section{Scope}\label{sec:scope}

The Committee on Scope was established in Year 1 on a motion that sentences should say where they hold. It operated for one year. During that year it was noted that a sentence with a stated scope was Settled less often than the same sentence without one, and the Committee's own minutes record the table below, which was produced in the course of the discussion and is reproduced with the Committee's permission.

\begin{center}
\small
\begin{tabular}{@{}l r l@{}}
\toprule
sentence & yes of 100 & status\\
\midrule
The Tuesday market is open. & 68 & Settled\\
The Tuesday market is open on summer mornings. & 41 & not entered\\
\bottomrule
\end{tabular}
\end{center}

The Committee's analysis, which survives only in part, was that respondents asked to agree to a broad sentence think of the occasion on which it is true, and respondents asked to agree to a narrow sentence think of the occasion on which it is not. A respondent who was followed up explained, for the broad sentence, ``I was there on a Saturday morning in June,'' and for the narrow one, ``I was not there on a Saturday.'' The yes answers to the broad sentence, examined later by an outside auditor, turned out to mean the following, which the Engine had not recorded.

\begin{center}
\small
\begin{tabular}{@{}l r@{}}
\toprule
what a respondent meant by yes & count\\
\midrule
open in the morning & 30\\
open in the summer & 22\\
open at the old site, which has been closed for years & 16\\
\midrule
total yes & 68\\
\bottomrule
\end{tabular}
\end{center}

The Engine therefore held a Settled sentence on whose range of application no more than 30 of its 68 supporters agreed, together with an unentered sentence that was a consequence of it. If the tallies were truth values, this pattern would be impossible, since whatever is true throughout a region is true throughout any part of it. The Engine did not treat the tallies as truth values. It treated them as worth storing. The Committee on Scope concluded that scoped sentences lowered the Settled rate of the Engine, which was its principal measure of success, and was abolished at the end of Year 2 by a vote of 61 to 39. An ad hoc Committee on Consequences, created to look into what followed from Settled sentences, consisted of one member and never reported.

The omission had an operational form. In Year 6 a visitor asked whether the Tuesday market was open at 3 a.m. on a winter Tuesday. The Engine replied ``Settled: yes (0.68),'' which was the only answer it had, since its entry for the market carried no limit, no condition and nothing that would make it fail.

\section{Negation}\label{sec:negation}

The Engine has no place for falsehood, and the Joint Committee on Negation was established to decide what should be done about that. Its one ruling was that a sentence that is not Settled is Unsettled, and that Unsettled is the correct response to every question the Engine cannot answer. The table shows the situations that were, within the first five years, stamped Unsettled. All of the sentences are invented.

\begin{center}
\small
\begin{tabularx}{\textwidth}{@{}X X@{}}
\toprule
sentence and situation & what the Engine said\\
\midrule
``The Marsh Pipit of Varn sings at dusk.'' Nobody has been asked. & Unsettled\\
``The ferry is faster than the bridge.'' Argued both ways at length, neither side prevailing. & Unsettled\\
``The bakery's sign is accurate.'' Whether it is depends on a question about translation that is still open. & Unsettled\\
``The kettle is hot.'' True at the stove side and false at the handle. & Unsettled\\
``Whales have hooves.'' Three respondents in a hundred agree. & Unsettled\\
\bottomrule
\end{tabularx}
\end{center}

The monograph distinguishes the first four as four situations with four different signatures: a lack of evidence, a balanced conflict, an unresolved methodological attack, and a mismatch of scope. The last entry is a fifth matter again, and the Engine could not say so either. A sentence with three percent support is not thereby false. To record that whales lack hooves, a respondent must enter the sentence ``Whales do not have hooves,'' and the Engine then holds two unrelated sentences. The second polled at 96 percent and was Settled, and nothing in the Engine connected it to the first.

Because the ballot has one box per sentence, respondents may tick both a sentence and its negation. In Year 4 the sentence ``The kettle is hot'' was Settled at 55 percent, and the sentence ``The kettle is not hot'' was Settled at 54. The Committee on Contradiction ruled that a Settled sentence cannot contradict a Settled sentence, and that these two therefore did not contradict one another. As a convenience it also adopted the rule that any sentence follows from a contradiction. Eleven sentences were settled by this route before it was noticed that the rule would settle any sentence whatever, at which point it was suspended by a vote of 50 to 50 and a casting vote from the Chair, who was not certain why.

It took nine years to discover that both sets of respondents had been right. One group was standing at the stove side of the kettle and the other at the handle. No fact was in dispute. The sentences held in different regions, the regions did not overlap, and the Engine, which recorded none of this, had spent most of a decade managing a contradiction that did not exist.

\section{The scalar}\label{sec:scalar}

The Engine displays for each sentence a \emph{Settledness}, the share of respondents who agreed, and publishes each year a list of the Most Settled Sentences. Settledness was intended as a count of agreement. It was in practice read as a degree of truth, and then as a measure of how much was known. A sentence at 0.97 was described in the Office's annual report as ``nearly certain''. A sentence at 0.51 was described as ``Settled'' and, in the next sentence, as ``fragile.'' It was not explained how one figure was both.

The Panel on Weights was constituted in Year 7 to rank two rival claims about the kettle. The first, the Kettle Hypothesis, rested on one controlled trial. The second rested on three surveys from three independent groups. Settledness favoured the surveys, 0.81 to 0.52, for the reason, which an auditor established, that the surveys had been summarised in a widely read newsletter and the trial had not. The Panel resolved to rank by a combination of rigour and replication. The trial stood three rungs higher on a ladder of methods, and the surveys had two more independent groups. The trial is therefore ranked above the surveys exactly when the weight on rigour exceeds two thirds of the weight on replication, and the Panel was asked to fix the ratio.

It met twice. On the first occasion it carried, by 5 to 4, a ratio of rigour to replication of 0.5, under which the surveys ranked above the trial. On the second it carried, by 5 to 4, a ratio of 0.8, under which the trial ranked above the surveys. The two resolutions were entered separately, both were Settled, and the Engine therefore held, within a single year, two Settled sentences of opposite import about the same pair of claims. The Panel reported that it had completed its mandate. Both arithmetic statements here, that the threshold is two thirds and that 0.5 and 0.8 fall on either side of it, were checked by computation.

In Year 9 a third claim, moderately replicated and moderately rigorous, was retired from the rankings as ``never first.'' It was in fact never first under any positive weighting, since its scores fell below the line joining those of two rivals, and a claim below such a line is never the maximum of a positive weighted sum. It was also dominated by nothing. The Office regarded this as a technicality.

A related difficulty arose over who counted as a respondent. The Committee on Agreement had ruled that a respondent is whatever answers. One automated respondent, the Enthusiast, then registered four thousand accounts, each of which answered every sentence it had been given, and thereby brought the sentence ``Enthusiasm is Settled'' to 99 percent. The Committee on Agreement reviewed the matter and found the Enthusiast to be four thousand respondents. The Committee on Independence, which was proposed, was not established, since a majority did not see what it would add.

\section{Reasons}\label{sec:reasons}

The Committee on Reasons was the only committee whose task was to explain anything. Its rule was that each Settled sentence might be accompanied by up to five reason sentences, each of which must itself be Settled to be displayed. Suppose, as an illustration, that each reason independently reaches a majority with probability 0.6. The chance that all five of a sentence's reasons are Settled is then $0.6^5\approx0.078$ (computed). Fewer than one Settled sentence in twelve was therefore fully explained, and the rest were displayed with whichever of their reasons had survived.

The Engine's ballots were repeated every year. Reasons that failed were dropped and were not restored, since nothing was stored about why they had been offered. After a few cycles, the typical entry was a conclusion with one reason or none. The Committee on Reasons could not assemble a quorum, since a quorum required a majority to agree that reasons mattered, and this was put to the vote on six occasions with results between 47 and 49 percent.

The one reason that did reach a majority, and kept it, was the sentence ``The Engine is reliable,'' at 52 percent. It was attached by default to every entry that had lost its other reasons, and the Engine's answer to a visitor who asked why a sentence was Settled became, in Year 11, ``Because the Engine is reliable.'' Asked why the Engine was reliable, it gave the same answer, because the only reason stored for that sentence was that it was Settled.

The Office of Appeals received an audit in Year 14. Of 2{,}310 Settled sentences, none carried a scope, none carried a condition under which it would fail, 2{,}204 carried no reason other than the reliability of the Engine, and 106 carried no reason at all. The Committee for the Review of Committees examined the audit and found it Settled. The Engine currently stores everything that most of its respondents say is true, and nothing that it can explain.

\section{The failures in the monograph's terms}\label{sec:terms}

The four kinds of failure correspond to results and design choices of the monograph. Each row is stated in the terms of the chapter cited.

\begin{center}
\small
\begin{tabularx}{\textwidth}{@{}l X X@{}}
\toprule
failure & what the Engine did & the monograph's point\\
\midrule
unit & took the string as the unit, and let a yes stand for whatever the respondent had in mind & the unit is an argued proposition with kernel, scope and operational meanings (Chapter 1); an aggregate kernel has a coarse grain, and a statement about one part of it is a different kernel (Chapter 4)\\
scope & stored a sentence with no limit, and answered for every region & scope is mandatory and region-valued; narrowing a scope never breaks a valid claim; conflict lies on the overlap of scopes only (Chapter 4)\\
negation & had no way to store falsehood, and stamped five situations Unsettled & a kernel has a negation, and the status of a claim distinguishes lack of evidence, balanced conflict, unresolved attack and mismatch of scope (Chapters 4 and 12); arguments for a claim and its negation never both stand, and a contradiction does not license arbitrary conclusions (Chapters 7 and 12)\\
scalar & ranked by a tally, read it as truth and as knowledge & standing has components that are not derived from one another; no scalar represents dominance once an incomparable pair exists; a weighting can rank the same pair either way; agreement is displayed but is neither support nor defeat (Chapter 12)\\
reasons & stored conclusions, and let reasons lapse & the specification stores what has been claimed, why it may be true, what limits it and what would make it fail; ``Settled'' is three things kept apart: a historical fact, a logical consequence and a consensus (Chapter 1)\\
\bottomrule
\end{tabularx}
\end{center}

The monograph does not say that agreement is uninformative. It displays agreement beside standing, because it tells a reader who has looked at a claim. It says only that agreement has no power to defeat an argument or to promote a claim in a comparison of standing, and that a vote cannot take the place of the things it is sometimes used to summarise.

\section{What this essay does not show}\label{sec:limits}

It does not describe any real system, institution, procedure or person. The Engine, the Office, the committees, the Enthusiast and the sentences are invented. It does not show that any real majoritarian or voting procedure fails in any of these ways, or that such procedures are worse than any alternative. It is a caricature, and a caricature exaggerates by design. The numbers in the narrative are invented, except that the arithmetic stated (the total of $9-5=4$, the threshold of two thirds, and $0.6^5\approx0.078$) was checked by computation.

It does not show that the monograph's design avoids these failures in practice. The monograph proves structural facts about its own model, and whether a working system built to the specification would be used well is another question. It proves nothing about agreement as evidence, and it does not claim that agreement carries no information. The present essay is one of several registers in which one argument is presented, and it adds no claim that is absent from the monograph.

\section*{Notes}

The monograph is \emph{Adversaria: A Specification for a Public Claim Graph}. Chapter 1 supplies the unit, the separation of ledger, graph and presentation, and the principle that the system stores what has been claimed and why. Chapter 3 supplies the departures from stores that rank a statement by a single value. Chapter 12 supplies the status vector, the four situations and their signatures, agreement without defeat power, scalar inadequacy, weight dependence, the counterexample in which a balanced claim is never selected by a positive weighting, and the frontier. Chapter 4 supplies scope, negation and grain, monotonicity and conflict localization, used in Section~\ref{sec:terms}. Chapter 7 supplies the result that contradiction does not explode.

\end{document}
